\documentclass[12pt,a4paper]{article}

\usepackage[utf8]{inputenc}
\usepackage[T1]{fontenc}
\usepackage[brazil]{babel}
\usepackage{geometry}
\usepackage{graphicx}
\usepackage{hyperref}
\usepackage{booktabs}
\usepackage{listings}
\usepackage{xcolor}
\usepackage{setspace}
\usepackage{indentfirst}
\usepackage{caption}

\geometry{left=3cm, right=2cm, top=3cm, bottom=2cm}
\onehalfspacing
\hypersetup{colorlinks=true, linkcolor=black, urlcolor=blue}

\lstdefinestyle{codigo}{
  basicstyle=\ttfamily\small,
  keywordstyle=\color{blue}\bfseries,
  commentstyle=\color{gray},
  stringstyle=\color{red!70!black},
  numbers=left,
  numberstyle=\tiny\color{gray},
  stepnumber=1,
  numbersep=8pt,
  frame=single,
  breaklines=true,
  tabsize=2,
  showstringspaces=false,
  literate={á}{{\'a}}1 {é}{{\'e}}1 {í}{{\'i}}1 {ó}{{\'o}}1 {ú}{{\'u}}1
           {ã}{{\~a}}1 {õ}{{\~o}}1 {ç}{{\c{c}}}1 {ê}{{\^e}}1 {ô}{{\^o}}1
           {Á}{{\'A}}1 {É}{{\'E}}1 {Í}{{\'I}}1 {Ó}{{\'O}}1 {Ú}{{\'U}}1
           {Ã}{{\~A}}1 {Õ}{{\~O}}1 {Ç}{{\c{C}}}1
}

\lstdefinestyle{gramatica}{
  basicstyle=\ttfamily\small,
  frame=single,
  breaklines=true,
  tabsize=2
}

\title{%
  \textbf{Relatório Técnico}\\[0.4cm]
  \large Implementação Parcial de um Compilador para Subconjunto da Linguagem C
}
\author{%
  Igor de Matos da Rosa\\[0.2cm]
  \small Matrícula: 20103930\\[0.6cm]
  \small Disciplina: DEC0004-09655 (20261) -- Compiladores\\
  \small Universidade Federal de Santa Catarina (UFSC)\\[0.3cm]
  \small Professor: Alison Roberto Panisson
}
\date{Araranguá, 2026}

\begin{document}

\maketitle
\thispagestyle{empty}
\newpage

\tableofcontents
\newpage

% ---------------------------------------------------------------------------
\section{Introdução}

Este relatório descreve a implementação parcial de um compilador desenvolvida como
Trabalho Final da disciplina DEC0004-09655 (20261) -- Compiladores, ministrada na Universidade
Federal de Santa Catarina. O objetivo da atividade consiste em integrar, em um único
artefato de software, as fases de \textbf{análise léxica}, \textbf{análise sintática}
e \textbf{análise semântica}, conforme especificado no enunciado da tarefa avaliativa.

A solução foi construída de forma incremental: a análise léxica constitui evolução
do Trabalho 1 (T1); a análise sintática, do Trabalho 2 (T2); e a análise semântica
e o relatório consolidado de execução foram incorporados nesta entrega final. A
linguagem-alvo é um \textbf{subconjunto restrito de C}, suficiente para demonstrar
as estruturas exigidas: declarações de variáveis, atribuições, expressões aritméticas,
estruturas condicionais (simples e aninhadas) e laços de repetição (\texttt{while} e
\texttt{for}).

A implementação utiliza a biblioteca \textbf{PLY} (\textit{Python Lex-Yacc}), ferramenta
didática amplamente empregada na disciplina para a construção de analisadores léxicos
e sintáticos a partir de expressões regulares e gramáticas livres de contexto.

% ---------------------------------------------------------------------------
\section{Fundamentação Teórica e Relação com o Conteúdo da Disciplina}

A organização deste trabalho segue a estrutura curricular disponibilizada no ambiente
virtual da disciplina, articulando teoria e prática nas unidades estudadas.

\subsection{Unidade 1 -- Introdução à Construção de Compiladores}

Conforme os materiais introdutórios da disciplina, um compilador transforma um
programa-fonte em um programa-alvo equivalente, percorrendo fases sequenciais de
análise e síntese. Neste trabalho, implementam-se as três primeiras fases da
\textbf{análise}:

\begin{enumerate}
  \item \textbf{Análise léxica} -- particionamento do fluxo de caracteres em tokens;
  \item \textbf{Análise sintática} -- verificação da estrutura conforme uma gramática;
  \item \textbf{Análise semântica} -- verificação de regras de contexto não expressas
        pela gramática.
\end{enumerate}

As fases de geração e otimização de código (Unidade 5) não foram implementadas,
pois extrapolam o escopo do trabalho final, que exige uma \textit{implementação
parcial}.

\subsection{Unidade 2 -- Análise Léxica}

A análise léxica realiza o reconhecimento de \textbf{lexemas} por meio de um
\textbf{autômato finito}, usualmente especificado por expressões regulares. No modelo
teórico, o analisador léxico remove comentários e espaços em branco, identifica
palavras reservadas, literais, identificadores e operadores, encaminhando tokens ao
analisador sintático.

Na implementação, cada categoria léxica é definida por uma função \texttt{t\_*} ou
por uma expressão regular associada a um terminal PLY. Palavras reservadas como
\texttt{int}, \texttt{if}, \texttt{while} e \texttt{for} são mapeadas para tokens
distintos de identificadores genéricos (\texttt{ID}). Operadores compostos
(\texttt{==}, \texttt{<=}, \texttt{>=}, \texttt{!=}, \texttt{++}, \texttt{--})
possuem regras declaradas \textbf{antes} dos operadores simples, evitando conflitos
de \textit{prefixo mais longo} (\textit{maximal munch}).

\subsection{Unidade 3 -- Análise Sintática}

A análise sintática verifica se a sequência de tokens pode ser derivada a partir da
gramática da linguagem. A disciplina aborda analisadores \textbf{top-down} (LL) e
\textbf{bottom-up} (LR). O PLY Yacc implementa um analisador do tipo \textbf{LALR(1)},
subclasse dos analisadores LR, adequada a gramáticas de linguagens de programação.

Os conceitos de \textbf{conjuntos First e Follow}, bem como a resolução de ambiguidades,
manifestam-se na gramática implementada. Destaca-se a ambiguidade clássica do
\texttt{if-then-else}, resolvida mediante declaração de precedência:

\begin{lstlisting}[style=gramatica]
precedence = (
    ('nonassoc', 'IFX'),
    ('nonassoc', 'ELSE'),
)
\end{lstlisting}

A regra condicional com \texttt{ELSE} recebe menor precedência que o \texttt{IF}
sem \texttt{ELSE}, associando o \texttt{else} ao \texttt{if} mais próximo --- regra
usual em linguagens como C.

A gramática reconhece, entre outras estruturas, o programa principal
\texttt{int main()}, blocos compostos, declarações, atribuições, expressões com
precedência aritmética, condicionais aninhados e laços \texttt{while}/\texttt{for}.

\subsection{Unidade 4 -- Análise Semântica}

A análise semântica verifica propriedades que a gramática não captura: existência
de identificadores, compatibilidade de tipos e consistência de escopos. A disciplina
distingue \textbf{atributos sintetizados} (calculados nos filhos e propagados ao pai)
e \textbf{atributos herdados} (propagados do pai aos filhos).

Nesta implementação, adota-se um esquema predominantemente \textbf{S-atribuído}:
cada nó de expressão sintetiza seu tipo em \texttt{p[0]}, permitindo verificações nas
produções de atribuição, declaração e estruturas de controle. A \textbf{tabela de
símbolos} materializa o conceito de \textbf{escopo léxico}, implementado como pilha de
tabelas hash: cada bloco \texttt{\{\}} abre um novo escopo e o encerramento do bloco
restaura o escopo anterior.

% ---------------------------------------------------------------------------
\section{Especificação do Trabalho}

O enunciado do Trabalho Final estabelece os seguintes requisitos:

\begin{itemize}
  \item análise sintática capaz de validar atribuições, operações aritméticas,
        condicionais (simples e aninhadas) e laços de repetição;
  \item análise semântica com pelo menos três ações semânticas;
  \item reconhecimento do programa-exemplo fornecido no enunciado;
  \item entrega de código-fonte, \texttt{README.txt} e relatório teórico-prático.
\end{itemize}

O programa de referência do enunciado é reproduzido abaixo:

\begin{lstlisting}[style=codigo, language=C]
int main()
{
    int a, b, resultado;
    a = 4;
    b = 6;
    if(a>= 10){
         resultado = a -  b;
    }else{
         resultado = a +  b;
    }
}
\end{lstlisting}

% ---------------------------------------------------------------------------
\section{Arquitetura da Implementação}

O projeto está organizado em módulos com responsabilidade única:

\begin{table}[h]
  \centering
  \caption{Estrutura de arquivos do projeto}
  \begin{tabular}{@{}ll@{}}
    \toprule
    \textbf{Arquivo} & \textbf{Responsabilidade} \\
    \midrule
    \texttt{compilador\_c.py}   & Analisador léxico, gramática e orquestração \\
    \texttt{tabela\_simbolos.py} & Gerenciamento de escopos e declarações \\
    \texttt{semantica.py}     & Ações semânticas e verificação de tipos \\
    \texttt{relatorio.py}     & Consolidação do relatório de execução \\
    \texttt{Testes/}           & Programas de validação \\
    \texttt{README.txt}         & Instruções de execução \\
    \bottomrule
  \end{tabular}
\end{table}

\subsection{Fluxo de Processamento}

O compilador parcial executa o seguinte pipeline ao receber um arquivo \texttt{.c}:

\begin{enumerate}
  \item \textbf{Fase léxica}: leitura do arquivo e geração da sequência de tokens;
  \item \textbf{Fase sintática}: aplicação da gramática LALR(1) sobre os tokens;
  \item \textbf{Fase semântica}: execução das ações associadas às produções durante
        o \textit{parsing} (esquema ad hoc com atributos sintetizados);
  \item \textbf{Relatório final}: apresentação consolidada dos resultados das três
        fases e veredito de validade do programa.
\end{enumerate}

A execução é parametrizável:

\begin{lstlisting}[style=codigo, language=bash]
python3 compilador_c.py [--lex | --sint] <arquivo.c>
\end{lstlisting}

% ---------------------------------------------------------------------------
\section{Detalhamento da Análise Léxica}

\subsection{Tokens Reconhecidos}

O analisador léxico reconhece palavras reservadas, identificadores, literais
(inteiros, ponto flutuante, caracteres e \textit{strings}), operadores aritméticos
e relacionais, delimitadores e comentários (linha e bloco). Espaços, tabulações e
retornos de carro são ignorados; quebras de linha atualizam o contador de linhas
para diagnósticos de erro.

\subsection{Relação com a Teoria}

Cada regra léxica corresponde a um padrão regular que pode ser representado por um
AFD. A ferramenta PLY constrói internamente o autômato a partir das expressões
fornecidas, materializando na prática o conteúdo da Unidade 2 sobre geradores de
analisadores léxicos e a transição de especificação declarativa para reconhecedor
executável.

% ---------------------------------------------------------------------------
\section{Detalhamento da Análise Sintática}

\subsection{Principais Produções Gramaticais}

A gramática é centrada no não terminal \texttt{programa}, derivando para a função
\texttt{main} e seu bloco. Destacam-se as produções:

\begin{lstlisting}[style=gramatica]
programa           : funcao_principal semicolon_opcional
funcao_principal   : INT MAIN LPAREN parametro_main RPAREN bloco
comando            : declaracao SEMICOLON
                   | atribuicao SEMICOLON
                   | comando_condicional
                   | comando_repeticao
                   | bloco
comando_condicional: IF LPAREN condicao RPAREN comando %prec IFX
                   | IF LPAREN condicao RPAREN comando ELSE comando
comando_repeticao  : WHILE LPAREN condicao RPAREN comando
                   | FOR LPAREN inicio_for SEMICOLON condicao_opcional
                     SEMICOLON atualizacao_for RPAREN comando
expressao          : expressao PLUS termo
                   | expressao MINUS termo
                   | termo
\end{lstlisting}

A hierarquia \texttt{expressao $\rightarrow$ termo $\rightarrow$ potencia $\rightarrow$ fator}
implementa precedência de operadores: adição/subtração, multiplicação/divisão e
potenciação, respectivamente.

\subsection{Relação com a Teoria}

O analisador LALR(1) construído pelo Yacc realiza \textit{shift} e \textit{reduce}
conforme a tabela de ações estudada na Unidade 3. Cada função \texttt{p\_*} representa
uma produção da gramática; a redução bem-sucedida confirma que o trecho analisado
pertence à linguagem definida.

% ---------------------------------------------------------------------------
\section{Detalhamento da Análise Semântica}

\subsection{Ações Semânticas Implementadas}

Foram implementadas \textbf{seis} ações semânticas, superando o mínimo de três
exigido pelo enunciado:

\begin{table}[h]
  \centering
  \caption{Ações semânticas implementadas}
  \begin{tabular}{@{}p{0.5\textwidth}p{0.38\textwidth}@{}}
    \toprule
    \textbf{Ação} & \textbf{Fundamento teórico} \\
    \midrule
    Inserção na tabela de símbolos & Escopo léxico e declaração \\
    Verificação de uso antes da declaração & Regra de contexto \\
    Detecção de redeclaração no mesmo escopo & Política de escopos aninhados \\
    Compatibilidade de tipos em atribuições & Checagem de tipos \\
    Tipagem de expressões aritméticas/relacionais & Atributos sintetizados \\
    Validação de condições em estruturas de controle & Restrição semântica de \texttt{if}/\texttt{while}/\texttt{for} \\
    \bottomrule
  \end{tabular}
\end{table}

\subsection{Tabela de Símbolos}

A classe \texttt{TabelaSimbolos} mantém uma pilha de dicionários. Ao encontrar
\texttt{\{}, um novo escopo é empilhado; ao encontrar \texttt{\}}, o escopo é
desempilhado. A operação de busca percorre os escopos do mais interno ao mais
externo, reproduzindo a regra de resolução de nomes de C.

\subsection{Verificação de Tipos}

Os tipos primitivos considerados são: \texttt{int}, \texttt{float}, \texttt{double},
\texttt{char} e \texttt{string}. Expressões aritméticas aplicam promoção numérica
(\textit{char} $\leq$ \textit{int} $\leq$ \textit{float} $\leq$ \textit{double}).
Operadores relacionais exigem operandos numéricos e produzem tipo \texttt{int},
analogamente à convenção de C.

% ---------------------------------------------------------------------------
\section{Testes e Validação}

Foram elaborados programas de teste em \texttt{Testes/}, classificados em válidos
(espera-se compilação sem erros semânticos) e inválidos (espera-se diagnóstico
semântico específico).

\begin{table}[h]
  \centering
  \caption{Casos de teste}
  \small
  \begin{tabular}{@{}llp{6.5cm}@{}}
    \toprule
    \textbf{Arquivo} & \textbf{Resultado} & \textbf{Objetivo} \\
    \midrule
    \texttt{ExemploEnunciado.c} & Válido & Programa do enunciado \\
    \texttt{Teste1.c}           & Válido & Condicionais aninhados \\
    \texttt{Teste2.c}           & Válido & Laço \texttt{while} \\
    \texttt{Teste3.c}           & Válido & \texttt{for}, \texttt{while} e \texttt{if} \\
    \texttt{Erro\_VariavelNaoDeclarada.c} & Inválido & Uso sem declaração \\
    \texttt{Erro\_Redeclaracao.c}         & Inválido & Redeclaração no escopo \\
    \texttt{Erro\_TipoIncompativel.c}     & Inválido & Atribuição \texttt{int} $\leftarrow$ \texttt{string} \\
    \texttt{Erro\_CondicaoInvalida.c}     & Inválido & Condição não numérica \\
    \bottomrule
  \end{tabular}
\end{table}

\subsection{Relatório de Execução}

Ao término da análise, o módulo \texttt{relatorio.py} imprime um relatório estruturado
contendo: (i)~lista de tokens com linha, tipo e valor; (ii)~estruturas sintáticas
reconhecidas; (iii)~tabela de símbolos e ações semânticas executadas; e (iv)~veredito
final (\textsc{Válido} ou \textsc{Inválido}). Esse relatório materializa, de forma
operacional, a integração das três fases analíticas estudadas na disciplina.

% ---------------------------------------------------------------------------
\section{Considerações Finais}

O trabalho consolidou os conhecimentos das Unidades 2, 3 e 4 da disciplina de
Compiladores, transformando conceitos teóricos --- autômatos finitos, gramáticas
livres de contexto, analisadores LR, tabelas de símbolos e esquemas de atributos ---
em um compilador parcial funcional.

As limitações do subconjunto implementado (ausência de geração de código, tipos
compostos, funções definidas pelo usuário e pré-processador) delimitam o escopo de
uma \textit{implementação parcial}, coerente com o enunciado. Como continuidade
natural do estudo, a Unidade 5 poderia estender o projeto com geração de código
intermediário ou código objeto.

% ---------------------------------------------------------------------------
\section*{Referências}

\begin{enumerate}
  \item PANISSON, A. R. \textit{Material didático da disciplina DEC0004-09655 (20261) -- Compiladores}. UFSC, 2026. Unidades 1 a 4 (Introdução, Análise Léxica, Análise Sintática e Análise Semântica).
  \item AHO, A. V.; LAM, M. S.; SETHI, R.; ULLMAN, J. D. \textit{Compiladores: Princípios, Técnicas e Ferramentas}. 2. ed. São Paulo: Pearson, 2008.
  \item BEAZLEY, D. M. \textit{PLY (Python Lex-Yacc)}. Disponível em: \url{https://www.dabeaz.com/ply/}. Acesso em: jun. 2026.
\end{enumerate}

\end{document}
